\subsection{Existence and uniqueness of p-rings}

PROPOSITION 4. --- Let C and C' be two p-rings such that \(l(\mathbf{C}) \geqslant l(\mathbf{C}')\) , and let \(\pi\) (resp. \(\pi'\) ) be the canonical homomorphism from C (resp. C') onto \(\kappa_{\mathbf{C}}\) (resp. \(\kappa_{\mathbf{C}'}\) ). Let \((x_{\lambda})_{\lambda \in \Lambda}\) (resp. \((x_{\lambda}')_{\lambda \in \Lambda}\) ) be a family of elements of C (resp. C') whose images under \(\pi\) (resp. \(\pi'\) ) form a p-basis of \(\kappa_{\mathbf{C}}\) (resp. \(\kappa_{\mathbf{C}'}\) ). Let v be an isomorphism of \(\kappa_{\mathbf{C}}\) onto \(\kappa_{\mathbf{C}'}\) such that \(v(\pi(x_{\lambda})) = \pi'(x_{\lambda}')\) for every \(\lambda \in \Lambda\) . There then exists a unique homomorphism u from C to C' such that \(v \circ \pi = \pi' \circ u\) and \(u(x_{\lambda}) = x_{\lambda}'\) for every \(\lambda \in \Lambda\) . It is surjective. If \(l(\mathbf{C}) = l(\mathbf{C}')\) , it is an isomorphism.

Let us prove the existence of u. Let A be the subring of C × C' formed by the pairs \((x, x')\) such that \(v(\pi(x)) = \pi'(x')\) . The map \((x, x') \mapsto \pi(x)\) is a surjective ring homomorphism of A onto \(\kappa_{\mathbb{C}}\) . Its kernel m, equal to \(p\mathbb{C} \times p\mathbb{C}'\) , is therefore a maximal ideal of A. The topological subspace A of \(\mathbb{C} \times \mathbb{C}'\) is closed in \(\mathbb{C} \times \mathbb{C}'\) , hence complete, and the topology induced on A by that of \(\mathbb{C} \times \mathbb{C}'\) is the m-adic topology. Consequently A is a separated and complete local ring with maximal ideal m (III, § 2, no. 13, prop. 19). For every \(\lambda \in \Lambda\) , one has \((x_{\lambda}, x_{\lambda}') \in A\) by hypothesis; if \(\xi_{\lambda}\) is the class of \((x_{\lambda}, x_{\lambda}')\) modulo m, the family \((\xi_{\lambda})_{\lambda \in \Lambda}\) is a \(p\) -basis of the field A/m. By th. 1 of no. 2, there exists a unique Cohen subring \(\mathbb{C}''\) of A containing \((x_{\lambda}, x_{\lambda}')\) for every \(\lambda \in \Lambda\) . We have \(l(\mathbb{C}'') = l(\mathbb{C}) \geqslant l(\mathbb{C}')\) . The restriction to \(\mathbb{C}''\) of the projection of \(\mathbb{C} \times \mathbb{C}'\) onto C is a homomorphism \(h: \mathbb{C}'' \to \mathbb{C}\) which induces an isomorphism of \(\kappa_{\mathbb{C}''}\) onto \(\kappa_{\mathbb{C}}\) . By prop. 2, c) of no. 2, h is an isomorphism of \(\mathbb{C}''\) onto C. One sees in the same way that the restriction \(h'\) to \(\mathbb{C}''\) of the projection of \(\mathbb{C} \times \mathbb{C}'\) onto C' is a surjective homomorphism of \(\mathbb{C}''\) into C'. Consequently, \(\mathbb{C}''\) is the graph of a surjective homomorphism \(u = h' \circ h^{-1}\) of C onto C', and we evidently have \(v \circ \pi = \pi' \circ u\) , \(u(x_{\lambda}) = x_{\lambda}'\) for every \(\lambda \in \Lambda\) . Moreover, if \(l(\mathbb{C}) = l(\mathbb{C}')\) , u is an isomorphism.

Let us prove the uniqueness of \(u\) . Let \(u_{1}\) be a homomorphism from C to C' such that \(v \circ \pi = \pi' \circ u_{1}\) and \(u_{1}(x_{\lambda}) = x_{\lambda}'\) for every \(\lambda \in \Lambda\) , and let \(\mathbf{C}_{1}\) be the graph of \(u_{1}\) . It is immediate that \(\mathbf{C}_{1}\) is a Cohen subring of A, containing \((x_{\lambda}, x_{\lambda}')\) for every \(\lambda \in \Lambda\) , whence \(\mathbf{C}_{1} = \mathbf{C}''\) (th. 1 of no. 2) and finally \(u_{1} = u\) .

PROPOSITION 5. --- Let k be a field of characteristic p, and let n be an integer ≥ 1, or + ∞. There exists a p-ring of length n whose residue field is isomorphic to k.

The ring \(\mathbf{W}(k)\) of Witt vectors with coefficients in \(k\) is a separated and complete integral local ring, whose residue field is isomorphic to \(k\) ( \(\S\) 1, no. 8, prop. 8), and we have \(p\) . \(1_{\mathbf{W}(k)} \neq 0\) (loc. cit., formula (52)). Let C be a Cohen subring of \(\mathbf{W}(k)\) (no. 2, th. 1). Then C is a \(p\) -ring of length + \(\infty\) whose residue field is isomorphic to \(k\) , and, if \(n\) is an integer \(\geqslant 1\) , the quotient \(\mathrm{C}/p^{n}\mathrm{C}\) is a \(p\) -ring of length \(n\) whose residue field is isomorphic to \(k\) .

Remarks. --- 1) Let \(n\) be an integer \(\geqslant 1\) and S be a \(p\) -basis of \(k\) . One can show that the subring of \(\mathbf{W}_{n}(k)\) generated by \(\mathbf{W}_{n}(k^{p^{n}})\) and by the elements \([\xi, 0, ..., 0]\) ( \(\xi \in S\) ), is a \(p\) -ring of length \(n\) whose residue field is isomorphic to \(k\) (cf. p. 72, exerc. 10).

\begin{enumerate}
\def\labelenumi{\arabic{enumi})}
\setcounter{enumi}{1}
\tightlist
\item
  The reader will find in the Appendix a proof of prop. 5 which uses neither the results of § 1, nor the existence theorem for Cohen subrings (no. 2, th. 1).
\end{enumerate}

COROLLARY. --- Let C be a p-ring of finite length n. There exists a p-ring C' of infinite length such that C is isomorphic to \(C'/p^{n}C'\).

By prop. 5, there exists a \(p\) -ring \(C'\) of infinite length such that \(\kappa_{C'}\) is isomorphic to \(\kappa_{C}\) . Then \(C'/p^{n}C'=C_{n}'\) is a \(p\) -ring of length \(n\) , and the field \(\kappa_{C_{n}'}\) is isomorphic to \(\kappa_{C'}\) , hence to \(\kappa_{C}\) . By prop. 4, the rings \(C\) and \(C_{n}'\) are therefore isomorphic.

